% ============================================================================
\section{Bayesian Networks (Car Tracking)}
\label{sec:bayesian}
% Source: assignments/06_bayesian/index.html
% ============================================================================

Set-up: our car's position $a_t$ is observed. The other car's location is $C_t$ (hidden). The microphone reading
is $D_t\sim\mathcal N(\norm{a_t-C_t},\sigma^2)$. Lower-case $p(\cdot)$ denotes local conditional distributions
of the Bayesian network; upper-case $\P(\cdot)$ denotes joint or posterior quantities obtained by inference.

% ----------------------------------------------------------------------------
\subsection{Emission probabilities}
The other car is stationary: $C_t=C$ for all $t$.

\problempart{1a}{Bayesian network}
\question{Draw the Bayesian network over $C,D_1,D_2,D_3,a_1,a_2,a_3$, shading the variables observed after $t=3$.}
\expect{A legible drawing with one arrow from each parent of each variable.}
\begin{solution}
Each reading $D_t$ depends on the (hidden) location $C$ and on our own position $a_t$, so both are parents of
$D_t$. After $t=3$ the readings and our positions are observed (shaded), and only $C$ is unknown.
\begin{center}
\begin{tikzpicture}[node distance=15mm and 18mm]
  \node[latent] (C) {$C$};
  \node[obs, below=of C] (D2) {$D_2$};
  \node[obs, left=of D2] (D1) {$D_1$};
  \node[obs, right=of D2] (D3) {$D_3$};
  \node[obs, below=of D1] (a1) {$a_1$};
  \node[obs, below=of D2] (a2) {$a_2$};
  \node[obs, below=of D3] (a3) {$a_3$};
  \foreach \t in {1,2,3} { \draw[edge] (C) -- (D\t); \draw[edge] (a\t) -- (D\t); }
\end{tikzpicture}
\end{center}
\end{solution}

\problempart{1b}{Joint distribution}
\question{Express $\P(C=c,D_1=d_1,D_2=d_2,D_3=d_3)$ using only local conditional distributions
($a_1,a_2,a_3$ are known with probability 1).}
\expect{A formula that uses only local conditional distributions.}
\begin{solution}
Multiplying the local conditional distributions of the network (the $a_t$ are fixed, observed values):
\[
\P(C=c,D_1=d_1,D_2=d_2,D_3=d_3)=p(c)\prod_{t=1}^{3}p(d_t\mid c,a_t)
=p(c)\prod_{t=1}^{3}p_{\mathcal N}\bigl(d_t;\ \norm{a_t-c},\ \sigma^2\bigr).
\]
\end{solution}

\problempart{1c}{\texttt{ExactInference.observe}}
\codepart{6.1c}
For every tile $c$ we multiply the current belief by $p(d_t\mid c)=p_{\mathcal N}(d_t;\norm{a_t-c},\sigma^2)$,
computed with \texttt{util.pdf(mean, std, value)} in that order, and then normalise. The grader awards 4/4, including
its test that swaps \texttt{util.pdf} for a function that is not symmetric in its arguments.

% ----------------------------------------------------------------------------
\subsection{Transition probabilities}
The other car moves according to $p(c_{t+1}\mid c_t)$, given in code as \texttt{self.transProb}.

\problempart{2a}{Bayesian network}
\question{Draw the Bayesian network over $C_1,C_2,C_3,D_1,D_2,D_3,a_1,a_2,a_3$, shading the variables observed after $t=3$.}
\expect{A legible drawing with one arrow from each parent of each variable.}
\begin{solution}
The locations now form a Markov chain $C_1\to C_2\to C_3$ (a hidden Markov model), and each reading depends on the
location at its own time step and on our position at that time.
\begin{center}
\begin{tikzpicture}[node distance=15mm and 18mm]
  \node[latent] (C1) {$C_1$};
  \node[latent, right=of C1] (C2) {$C_2$};
  \node[latent, right=of C2] (C3) {$C_3$};
  \node[obs, below=of C1] (D1) {$D_1$};
  \node[obs, below=of C2] (D2) {$D_2$};
  \node[obs, below=of C3] (D3) {$D_3$};
  \node[obs, below=of D1] (a1) {$a_1$};
  \node[obs, below=of D2] (a2) {$a_2$};
  \node[obs, below=of D3] (a3) {$a_3$};
  \draw[edge] (C1) -- (C2); \draw[edge] (C2) -- (C3);
  \foreach \t in {1,2,3} { \draw[edge] (C\t) -- (D\t); \draw[edge] (a\t) -- (D\t); }
\end{tikzpicture}
\end{center}
\end{solution}

\problempart{2b}{Joint distribution}
\question{Express $\P(C_1=c_1,C_2=c_2,C_3=c_3,D_1=d_1,D_2=d_2,D_3=d_3)$ using only local conditional distributions.}
\expect{A formula that uses only local conditional distributions.}
\begin{solution}
\[
\P(c_1,c_2,c_3,d_1,d_2,d_3)=p(c_1)\,p(c_2\mid c_1)\,p(c_3\mid c_2)\,\prod_{t=1}^{3}p(d_t\mid c_t,a_t),
\]
\[
p(d_t\mid c_t,a_t)=p_{\mathcal N}\bigl(d_t;\ \norm{a_t-c_t},\ \sigma^2\bigr).
\]
\end{solution}

\problempart{2c}{\texttt{ExactInference.elapseTime}}
\codepart{6.2c}
We build the new belief in a separate grid of zeros. For every pair $(c_t,c_{t+1})$ in the sparse transition table we
add $\P(C_t=c_t\mid d_{1:t})\,p(c_{t+1}\mid c_t)$, and then normalise. Every term therefore uses only the old
belief, and the work is proportional to the number of transitions rather than (number of tiles)$^2$. The grader
awards 2/2, and the hidden stress tests run in under $0.04$\,s (limit 5\,s).

% ----------------------------------------------------------------------------
\subsection{Which car is it?}
With $K$ cars, we observe only $\mathbf E_t=[E_{t1},\dots,E_{tK}]$, a uniformly random permutation of the true
noisy distances $\mathbf D_t=[D_{t1},\dots,D_{tK}]$, so a reading cannot be attributed to a specific car. All
distances are distinct.

\begin{figure}[H]\centering
  \includegraphics[width=0.45\linewidth]{06_bayesian/Q4_diag.jpeg}
  \caption{Flow of information for $K=2$ (from the handout; not a Bayesian network).}
\end{figure}

\problempart{3a}{Posterior for $K=2$, $T=1$}
\question{Write $p(C_{11}=c_{11},C_{12}=c_{12}\mid\mathbf E_1=\mathbf e_1)$ in terms of the Gaussian pdf
$p_{\mathcal N}(v;\mu,\sigma^2)$ and the priors $p(c_{11}),p(c_{12})$, without $D_{11},D_{12}$.
(For $K=1$: $p(c_{11}\mid\mathbf e_1)\propto p(c_{11})\,p_{\mathcal N}(e_{11};\norm{a_1-c_{11}},\sigma^2)$.)}
\expect{The expression (proportionality is fine) and the steps taken to derive it.}
\begin{solution}
Write $\mu_i=\norm{a_1-c_{1i}}$ for the true distance to car $i$ and $\pi$ for the random permutation, so that
$\mathbf E_1=(D_{1\pi(1)},D_{1\pi(2)})$ and each of the two orderings has probability $\tfrac12$.
\begin{enumerate}[nosep]
  \item Bayes' rule: $p(c_{11},c_{12}\mid\mathbf e_1)\propto p(c_{11})\,p(c_{12})\,p(\mathbf e_1\mid c_{11},c_{12})$
        (the cars' locations are independent a priori).
  \item Condition on the permutation and marginalise it out:
        $p(\mathbf e_1\mid c_{11},c_{12})=\sum_{\pi}p(\pi)\,p(\mathbf e_1\mid\pi,c_{11},c_{12})$.
  \item Given $\pi$, $e_{1j}$ is the reading of car $\pi(j)$, and $D_{11},D_{12}$ are independent given the locations, so
        $p(\mathbf e_1\mid\pi,c_{11},c_{12})=\prod_j p_{\mathcal N}(e_{1j};\mu_{\pi(j)},\sigma^2)$. This integrates out
        $D_{11},D_{12}$, because $\mathbf E_1$ is a deterministic rearrangement of them once $\pi$ is fixed.
\end{enumerate}
Therefore
\[
p(c_{11},c_{12}\mid\mathbf e_1)\ \propto\ p(c_{11})\,p(c_{12})\Bigl[
p_{\mathcal N}(e_{11};\mu_1,\sigma^2)\,p_{\mathcal N}(e_{12};\mu_2,\sigma^2)
+p_{\mathcal N}(e_{11};\mu_2,\sigma^2)\,p_{\mathcal N}(e_{12};\mu_1,\sigma^2)\Bigr],
\]
with $\mu_i=\norm{a_1-c_{1i}}$ (the common factor $\tfrac12$ is absorbed into the proportionality). Compared with
$K=1$, the likelihood is a sum over the two ways of matching readings to cars.
\end{solution}

\problempart{3b}{At least $K!$ maximisers}
\question{If the prior $p(c_{1i})$ is the same for every car, show that at least $K!$ assignments
$(c_{11},\dots,c_{1K})$ attain the maximum of $p(C_{11}=c_{11},\dots,C_{1K}=c_{1K}\mid\mathbf E_1=\mathbf e_1)$
(assume the maximising locations are distinct).}
\expect{A short mathematical argument or a concise explanation.}
\begin{solution}
Generalising 3a,
\[
p(c_{11},\dots,c_{1K}\mid\mathbf e_1)\propto\Bigl(\prod_{i=1}^K p(c_{1i})\Bigr)\,
\frac1{K!}\sum_{\pi\in S_K}\prod_{j=1}^K p_{\mathcal N}\bigl(e_{1j};\norm{a_1-c_{1\pi(j)}},\sigma^2\bigr).
\]
Let $\tau$ be any permutation of the car labels and consider the relabelled assignment
$(c_{1\tau(1)},\dots,c_{1\tau(K)})$. The prior product is unchanged because all cars share the same prior and the
product is commutative. The likelihood is also unchanged, because composing with $\tau$ maps the set of all
permutations $S_K$ onto itself ($\pi\mapsto\tau\circ\pi$), so the sum contains exactly the same terms. Hence every
relabelling has the same posterior probability. If $(c^\star_{11},\dots,c^\star_{1K})$ is a maximiser with
distinct locations, its $K!$ relabellings are $K!$ \emph{distinct} assignments that all attain the maximum, so
there are at least $K!$ maximisers. Intuitively, the readings carry no labels and the cars are exchangeable a priori,
so the data can tell us \emph{where} the cars are but never \emph{which car is which}.
\end{solution}

% ----------------------------------------------------------------------------
\subsection{Ethics in advanced technologies}

\problempart{4a}{Camera data misuse (dual use)}
\question{Self-driving cars often use cameras, which capture faces, licence plates and more. How might camera
data or its readings be misused, and what measures could pre-empt the misuse (or why are there none)?}
\expect{4--6 sentences.}
\begin{solution}
A fleet of camera-equipped cars is in effect a moving, city-wide surveillance network. Combining the time-stamped
licence plates and faces it records lets the operator, a hacker, or a government that subpoenas the data
reconstruct where specific people live, work, worship or protest, and track them over time without their consent.
The same tracking output could be repurposed for targeted advertising, insurance pricing, stalking, or locating
members of a particular community. Measures that pre-empt this: process images on the vehicle and keep only the
abstract output the driving stack needs (bounding boxes, velocities), never the raw frames; blur faces and plates
at capture time; delete what must be kept after a short, fixed retention period; encrypt stored data and log every
access; and set legal and contractual limits on sharing data with third parties. These measures reduce, but cannot
remove, the dual-use risk: whoever controls the software can switch them off, so independent audits and
regulation are needed as well.
\end{solution}

\problempart{4b}{Sensor deception}
\codepart{6.4b}
Observed distance after the attack:
$D'_t=\frac{1}{1+\text{skewness}^2}\,D_t+\sqrt{\frac{2}{1+\text{skewness}^2}}$, implemented in
\texttt{ExactInferenceWithSensorDeception.observe}, which transforms the reading and then calls the ordinary
emission update. With the default skewness $0.5$ every reading is shrunk to $80\%$ of its value plus about
$1.26$, so the tracker systematically believes the other cars are \emph{closer} than they really are. The grader
awards 3/3.

\problempart{4c}{Guarding against sensor deception}
\question{How could you guard against sensor deception if the type of attack is known in advance? (Hint: sensor redundancy.)}
\expect{1--2 sentences.}
\begin{solution}
Use redundant, physically different sensors, such as a microphone, lidar, radar and cameras, and fuse them, so
that a reading that disagrees with the consensus is flagged and down-weighted. An attacker would then have to fool
several independent sensing modalities at once. Because the form of the attack is known, we can also invert it,
$D_t=(1+s^2)\bigl(D'_t-\sqrt{2/(1+s^2)}\bigr)$, or test incoming readings for that skew pattern before using them.
\end{solution}

\problempart{4d}{Particle filtering vs.\ exact inference}
\question{Run \texttt{python drive.py -a -d -k 3 -i particleFilter}. How does the simulation differ from exact inference, and why?}
\expect{1--2 sentences.}
\begin{solution}
With the particle filter, the belief cloud around each car is patchier and noisier than with exact inference. It
jumps between frames, leaves gaps in plausible areas, and can lock onto the wrong spot for a while. The reason is
that it represents the belief with only 200 sampled particles that are resampled at every step: tiles without a
particle get probability exactly zero, and sampling noise moves the cloud around, whereas exact inference keeps a
probability for every tile. We measured this without the GUI (\texttt{scripts/car\_tracking\_comparison.py}, Lombard map,
20 runs of 40 steps). The most likely tile is on average 78.9 map units from the true car for the particle filter versus 65.0
for exact inference, and the particle belief puts more than $1\%$ probability on fewer tiles (7.9 vs.\ 9.7). On a
map this small (253 tiles for 200 particles) the particle filter was not faster ($3.1$ vs.\ $1.6$\,ms per step);
its advantage only appears when the number of tiles is much larger than the number of particles.
\end{solution}

\problempart{4e}{Privacy--accuracy trade-off}
\question{Particle filtering can run on-device (more private and secure) but is less accurate than cloud-based
exact inference. How would you weigh this trade-off?}
\expect{3--5 sentences.}
\begin{solution}
We would treat safety as a hard constraint and privacy as the objective inside it, not trade one off freely against
the other. First, measure whether the on-device particle filter is accurate \emph{enough}: count near-misses, collisions
and false alarms in simulation and on test tracks, and raise the number of particles (which costs on-device compute,
not privacy) until the safety margin matches the cloud system. If that is achievable, on-device processing is
preferable: it removes a single point of failure (the network link), lowers the attack surface, and avoids collecting
location data about everyone on the road. If it is not, a hybrid design can send only minimal, anonymised summaries to the
cloud when needed. In any case the decision should be made with regulators and the affected public, not by the
manufacturer alone, because the people whose data is at stake are mostly not the car's owners.
\end{solution}

\problempart{4f}{Adversarial attack on camera-based training}
\question{A camera-based autonomous vehicle is further trained on customers' anonymised camera data. Name an
adversarial attack, its potential impact on the vehicle, and a mitigation.}
\expect{3--5 sentences.}
\begin{solution}
\textbf{Attack:} data poisoning with a backdoor. Because the car keeps learning from customer footage, an attacker who
controls some vehicles, or simply drives in front of many of them, can repeatedly show a trigger pattern, such as a
specific sticker on stop signs, in situations that are labelled ``no stop needed''. \textbf{Impact:} the retrained
model learns that the sticker means ``ignore this sign'', so it behaves normally in all tests but runs stop signs
carrying the sticker, which could be exploited against a particular location or community. \textbf{Mitigation:} do
not feed customer data straight into training. Instead, filter it for anomalies and outliers, limit how much any
single vehicle or location can contribute, validate every new model on held-out safety-critical scenarios including
known trigger patterns, and keep a rule-based safety layer (for example, map-based stop-sign knowledge) that the
learned model cannot override.
\end{solution}
